\documentclass{report}
\usepackage{listings}
\usepackage{graphicx}

\begin{document}
\title{labwork 8}
\maketitle
\section{RGB to HSV}
The goal of the exercise is to turn an image from RGB format to HSV and then to turn it back to RGB.
\subsection{Preparation}
First, we normalize the RGB values from [0, 255] to [0, 1]. Then, we take the max, the min and the delta between min and max.
\begin{lstlisting}
    r = src[y, x, 0] / 255
    g = src[y, x, 1] / 255
    b = src[y, x, 2] / 255
    maxV = max(r, g, b)
    minV = min(r, g, b)
    delta = maxV - minV
\end{lstlisting}
\subsection{Conversion}
To turn our RGB image to HSV we apply the following formula:
\begin{figure}[ht]
    \centering
    \includegraphics[width=0.75\linewidth]{formula.png}
    \label{fig:placeholder}
\end{figure}
\newpage
Which gives in python:
\begin{lstlisting}
    if (maxV == 0):
        dst[1, y, x] = maxV
    else:
        dst[1, y, x] = delta / maxV

    dst[2, y, x] = maxV
    if delta == 0:
        dst[0, y, x] = 0
    elif maxV == r:
        dst[0, y, x] = 60 * (((g - b) / delta) % 6)
    elif maxV == g:
        dst[0, y, x] = 60 * ((b - r) / delta + 2)
    else:
        dst[0, y, x] = 60 * ((r - g) / delta + 4)
\end{lstlisting}
\section{HST to RGB}
\subsection{Preparation}
To turn the HSV image back to RGB we simply implement the given formulas:
\begin{lstlisting}
    d = h / 60
    f = d - int(d)
    l = v * (1 - s)
    m = v * (1 - f * s)
    n = v * (1 - (1 - f) * s)
\end{lstlisting}
\subsection{Conversion}
Then we compute the right result depending on the value of H going from 0 to 360.
\begin{lstlisting}
    if (0 <= h < 60):
        dst[y, x, 0] = round(v * 255)
        dst[y, x, 1] = round(n * 255)
        dst[y, x, 2] = round(l * 255)
    elif (60 <= h < 120):
        dst[y, x, 0] = round(m * 255)
        dst[y, x, 1] = round(v * 255)
        dst[y, x, 2] = round(l * 255)
\end{lstlisting}
\end{document}